% ch06_schauder.tex — Schauder Fixed Point Theorem and Variants
\chapter{Schauder Fixed Point Theorem and Variants}\label{ch:schauder}

\section{Motivation}

Brouwer's theorem is a gem of finite-dimensional topology, but it fails spectacularly in infinite dimensions: the closed unit ball of an infinite-dimensional Banach space is not compact, and continuous maps without fixed points exist. Juliusz Schauder, a Polish mathematician of the Lw\'ow school, overcame this obstacle in 1930 by replacing compactness of the set with \emph{compactness of the operator}: if $T$ maps a closed bounded convex set into a relatively compact subset of itself, then $T$ has a fixed point. This result is the pillar of existence methods for nonlinear differential and integral equations.

Brouwer's theorem (Chapter~5) applies only in finite dimensions. In infinite
dimensions, the closed unit ball is not compact, and continuous maps without
fixed points exist. Schauder resolved this problem by replacing compactness of
the set with compactness of the operator.

\section{Compact operators}

\begin{definition}[Compact operator]
Let $E$ be a Banach space and $C \subset E$. A continuous operator $T : C \to E$
is \emph{compact} if $T(B)$ is relatively compact (i.e., $\overline{T(B)}$ is
compact) for every bounded subset $B \subset C$.
\end{definition}

\begin{definition}[Completely continuous operator]
An operator $T : C \to E$ is \emph{completely continuous} if it is continuous and
maps every bounded set to a relatively compact set.
\end{definition}

\begin{remark}
In a Banach space, every completely continuous operator is compact. The converse
holds when the domain is bounded.
\end{remark}

\begin{example}[Fredholm integral operator]
The operator $(Tu)(t) = \int_0^1 K(t,s) u(s) \, ds$ on $C([0,1])$ is compact
whenever $K \in C([0,1]^2)$. This is a consequence of the Arzel\`a--Ascoli theorem.
\end{example}

\section{Schauder's theorem}

\begin{theorem}[Schauder, 1930]\label{thm:schauder}
Let $E$ be a Banach space, $C \subset E$ a closed, bounded, convex subset, and
$T : C \to C$ a compact operator. Then $T$ has at least one fixed point.
\end{theorem}

The proof relies on the following approximation lemma.

\begin{lemma}[Schauder approximation]\label{lem:schauder-approx}
Let $K$ be a compact subset of a Banach space $E$ and $\varepsilon > 0$. There
exists a finite-dimensional subspace $F \subset E$ and a continuous map
$P_\varepsilon : K \to \mathrm{co}(K) \cap F$ such that
$\norm{P_\varepsilon(x) - x} < \varepsilon$ for all $x \in K$.
\end{lemma}

\begin{proof}
By compactness, there exist $x_1, \ldots, x_n \in K$ such that
$K \subset \bigcup_{i=1}^n B(x_i, \varepsilon)$. Define
$\mu_i(x) = \max(0, \varepsilon - \norm{x - x_i})$ and
\[
    P_\varepsilon(x) = \frac{\sum_{i=1}^n \mu_i(x) \, x_i}{\sum_{i=1}^n \mu_i(x)}.
\]
This map is well defined (the denominator is positive since $x$ is within
distance $< \varepsilon$ of at least one $x_i$), continuous, and
$P_\varepsilon(x) \in \mathrm{co}(\{x_1, \ldots, x_n\})$. Moreover,
\[
    \norm{P_\varepsilon(x) - x} = \norm{\frac{\sum \mu_i(x)(x_i - x)}{\sum \mu_i(x)}}
    \leq \frac{\sum \mu_i(x) \norm{x_i - x}}{\sum \mu_i(x)} < \varepsilon
\]
since $\mu_i(x) > 0$ implies $\norm{x_i - x} < \varepsilon$.
\end{proof}

\begin{proof}[Proof of Theorem~\ref{thm:schauder}]
Set $K = \overline{T(C)}$, which is compact. For each $n \geq 1$, let
$P_n = P_{1/n}$ be the Schauder approximation from Lemma~\ref{lem:schauder-approx}.
The map $P_n \circ T : C \to \mathrm{co}(K) \cap F_n$ is continuous with values
in a finite-dimensional space.

Set $C_n = C \cap F_n$, which is convex, closed, and bounded in $F_n$. The map
$P_n \circ T|_{C_n} : C_n \to C_n$ is continuous (adjusting the domain if
necessary to ensure invariance). By Brouwer's theorem, there exists
$x_n \in C_n$ with $P_n(T(x_n)) = x_n$.

We have $\norm{x_n - T(x_n)} = \norm{P_n(T(x_n)) - T(x_n)} < 1/n$. By
compactness of $T$, the sequence $(T(x_n))$ has a convergent subsequence:
$T(x_{n_k}) \to y$. Then $x_{n_k} \to y$ as well, and by continuity of $T$:
$T(y) = \lim T(x_{n_k}) = y$. Since $C$ is closed, $y \in C$.
\end{proof}

\begin{keyformulas}[title=\textbf{Schauder's theorem}]
$C$ closed bounded convex in a Banach space, $T : C \to C$ compact $\implies$
$\exists\, x^* \in C : T(x^*) = x^*$.
\end{keyformulas}

\section{Schauder--Tychonoff version}

\begin{theorem}[Schauder--Tychonoff]
Let $E$ be a separated locally convex space, $C \subset E$ a nonempty compact
convex subset, and $T : C \to C$ continuous. Then $T$ has a fixed point.
\end{theorem}

\begin{remark}
This theorem generalizes both Brouwer (finite dimension) and Schauder (Banach
space). The local convexity hypothesis on the ambient space is essential.
\end{remark}

\section{Leray--Schauder theorem}

\begin{definition}[Leray--Schauder degree]
Let $E$ be a Banach space, $\Omega \subset E$ a bounded open set, and
$T : \overline{\Omega} \to E$ a compact operator with $x \neq T(x)$ for all
$x \in \partial\Omega$. The \emph{Leray--Schauder degree}
$\deg_{LS}(\mathrm{Id} - T, \Omega, 0) \in \Z$ is defined by
finite-dimensional approximation.
\end{definition}

\begin{theorem}[Leray--Schauder, 1934]\label{thm:leray-schauder}
Let $E$ be a Banach space, $\Omega \subset E$ a bounded open set containing $0$,
and $T : \overline{\Omega} \to E$ a compact operator. Suppose that
\[
    x \neq \lambda T(x) \quad \forall x \in \partial\Omega, \; \forall \lambda \in [0, 1].
\]
Then $T$ has a fixed point in $\Omega$.
\end{theorem}

\begin{proof}[Proof sketch]
The homotopy $H(\lambda, x) = x - \lambda T(x)$ satisfies $H(\lambda, x) \neq 0$
for $x \in \partial\Omega$ and $\lambda \in [0,1]$. By homotopy invariance of
the degree:
\[
    \deg_{LS}(\mathrm{Id} - T, \Omega, 0) = \deg_{LS}(\mathrm{Id}, \Omega, 0) = 1.
\]
Since the degree is nonzero, the equation $x - T(x) = 0$ has a solution.
\end{proof}

\begin{corollary}[Leray--Schauder alternative]
Let $E$ be a Banach space and $T : E \to E$ a completely continuous operator.
Then at least one of the following holds:
\begin{enumerate}
    \item The equation $x = T(x)$ has a solution;
    \item The set $\{x \in E : x = \lambda T(x) \text{ for some } \lambda \in (0,1)\}$
          is unbounded.
\end{enumerate}
\end{corollary}

\section{Measure of noncompactness and Darbo's theorem}

\begin{definition}[Kuratowski measure of noncompactness]
Let $A$ be a bounded subset of a metric space. The \emph{Kuratowski measure of
noncompactness} is
\[
    \alpha(A) = \inf\left\{ \varepsilon > 0 : A \subset \bigcup_{i=1}^n A_i
    \text{ with } \mathrm{diam}(A_i) \leq \varepsilon \right\}.
\]
\end{definition}

\begin{proposition}[Properties of $\alpha$]
For bounded subsets $A, B$ of a Banach space:
\begin{enumerate}
    \item $\alpha(A) = 0 \Leftrightarrow \overline{A}$ is compact;
    \item $A \subset B \Rightarrow \alpha(A) \leq \alpha(B)$;
    \item $\alpha(A \cup B) = \max(\alpha(A), \alpha(B))$;
    \item $\alpha(\overline{\mathrm{co}}(A)) = \alpha(A)$;
    \item $\alpha(A + B) \leq \alpha(A) + \alpha(B)$;
    \item $\alpha(\lambda A) = \abs{\lambda} \alpha(A)$.
\end{enumerate}
\end{proposition}

\begin{definition}[Condensing map]
A continuous operator $T : C \to C$ (where $C$ is closed bounded convex) is
\emph{$\alpha$-condensing} if there exists $k \in [0, 1)$ such that
$\alpha(T(A)) \leq k \, \alpha(A)$ for every bounded $A \subset C$.
\end{definition}

\begin{theorem}[Darbo, 1955]\label{thm:darbo}
Let $E$ be a Banach space, $C \subset E$ a nonempty closed bounded convex set,
and $T : C \to C$ an $\alpha$-condensing operator. Then $T$ has a fixed point.
\end{theorem}

\begin{proof}
Define recursively $C_0 = C$ and $C_{n+1} = \overline{\mathrm{co}}(T(C_n))$.
The sequence $(C_n)$ is decreasing ($C_{n+1} \subset C_n$) since
$T(C_n) \subset C_n$ implies
$C_{n+1} = \overline{\mathrm{co}}(T(C_n)) \subset \overline{\mathrm{co}}(C_n) = C_n$.

Set $\alpha_n = \alpha(C_n)$. Then
$\alpha_{n+1} = \alpha(\overline{\mathrm{co}}(T(C_n))) = \alpha(T(C_n)) \leq k \alpha_n$.
Hence $\alpha_n \leq k^n \alpha_0 \to 0$.

Set $C_\infty = \bigcap_{n \geq 0} C_n$. As a decreasing intersection of closed
sets in a complete space with $\alpha(C_n) \to 0$, the set $C_\infty$ is
nonempty and compact (by generalized Cantor). Moreover, $C_\infty$ is convex.
We have $T(C_\infty) \subset T(C_n) \subset C_{n+1}$ for all $n$, so
$T(C_\infty) \subset C_\infty$.

By Schauder's theorem applied to $T : C_\infty \to C_\infty$ (continuous on a
compact convex set), $T$ has a fixed point.
\end{proof}

\begin{remark}
Schauder's theorem is a special case of Darbo's: if $T$ is compact, then
$\alpha(T(A)) = 0$ for every bounded $A$, so $T$ is $\alpha$-condensing with
$k = 0$.
\end{remark}

\section{Sadovskii's theorem}

\begin{theorem}[Sadovskii, 1967]
Let $C$ be a closed, bounded, convex subset of a Banach space, and $T : C \to C$
continuous with $\alpha(T(A)) < \alpha(A)$ for every $A \subset C$ with
$\alpha(A) > 0$. Then $T$ has a fixed point.
\end{theorem}

\begin{remark}
Sadovskii generalizes Darbo by replacing the condition
$\alpha(T(A)) \leq k\alpha(A)$ ($k < 1$) with $\alpha(T(A)) < \alpha(A)$. The
proof uses a Zorn's lemma argument.
\end{remark}

\section{Krasnoselskii's theorem}

\begin{theorem}[Krasnoselskii, 1958]\label{thm:krasnoselskii}
Let $C$ be a closed, bounded, convex subset of a Banach space $E$. Let
$A : C \to E$ be a compact operator and $B : C \to E$ a contraction with
constant $k < 1$. If $A(x) + B(y) \in C$ for all $x, y \in C$, then the
equation $x = A(x) + B(x)$ has a solution in $C$.
\end{theorem}

\begin{proof}
For each $x \in C$, the map $y \mapsto A(x) + B(y)$ is a contraction from $C$
to $C$ (by hypothesis). By Banach's theorem, it has a unique fixed point $y(x)$:
$y(x) = A(x) + B(y(x))$. The map $x \mapsto y(x)$ is continuous (by continuous
dependence of the fixed point on the parameter) and compact (since
$y(x) = A(x) + B(y(x))$ and $A$ is compact). By Schauder's theorem, $y(\cdot)$
has a fixed point $x^* = y(x^*)$, whence $x^* = A(x^*) + B(x^*)$.
\end{proof}

\begin{remark}
Krasnoselskii's theorem is particularly useful for integral equations of the
form $x = Ax + Bx$ where $A$ is an integral operator (compact) and $B$ is a
contraction.
\end{remark}

\section{Applications}

\begin{example}[Nonlinear integral equation]
Consider the equation
\[
    u(t) = g(t) + \int_0^1 K(t,s) \, h(s, u(s)) \, ds, \quad t \in [0,1]
\]
where $g \in C([0,1])$, $K \in C([0,1]^2)$, and $h : [0,1] \times \R \to \R$ is
continuous and bounded. The operator
$(Tu)(t) = g(t) + \int_0^1 K(t,s) h(s, u(s)) \, ds$ is compact on the ball
$\overline{B}(0, R) \subset C([0,1])$ for $R$ large enough, and
$T(\overline{B}(0,R)) \subset \overline{B}(0,R)$ provided
$R \geq \norm{g}_\infty + \norm{K}_\infty \norm{h}_\infty$. By Schauder's
theorem, the equation has a solution.
\end{example}

\section{Exercises}

\begin{exercise}
Show that the Volterra operator $(Tu)(t) = \int_0^t K(t,s) u(s) \, ds$ is
compact on $C([0,1])$ when $K \in C([0,1]^2)$.
\end{exercise}

\begin{exercise}
Verify that $\alpha(\overline{B}(0, r)) = 2r$ in an infinite-dimensional Banach
space, and $\alpha(\overline{B}(0, r)) = 0$ in finite dimensions.
\end{exercise}

\begin{exercise}
Use Schauder's theorem to prove existence of a solution to the Cauchy--Peano
problem: $y' = f(t,y)$, $y(0) = y_0$, where $f$ is continuous (but not
necessarily Lipschitz).
\end{exercise}

\begin{exercise}
Show that Darbo's theorem implies Schauder's theorem.
\end{exercise}

\begin{exercise}
Apply Krasnoselskii's theorem to show existence of a solution to the equation
$u(t) = \frac{1}{3} u(t) + \int_0^1 K(t,s) g(u(s)) \, ds$ in $C([0,1])$,
under suitable hypotheses on $K$ and $g$.
\end{exercise}

\begin{exercise}
Show that in a reflexive Banach space, the closed unit ball is weakly compact.
Deduce a version of Schauder's theorem for the weak topology.
\end{exercise}
